
\subsection{Evaluation Points in a MovieGen}

MovieGen contains multiple components that can be used to generate a movie. For video output models, we try to set up a consistent evaluation criteria across all models. In Figure~\ref{evaluation_breakdown}, we show the different main evaluation points we have in MovieGen.

\begin{figure}[ht!]
\includegraphics[width=0.97\linewidth]{figures/moviegen_evaluation_points.pdf}
\caption{Component Evaluations in MovieGen. \fixme{Sam -- please fix this. English is broken here.}}
\label{img:eval}
\label{evaluation_breakdown}
\end{figure}


\subsection{Evaluating Image Outputs}

\subsubsection{Benchmark Generation}
Graph Representation

\begin{figure}[ht!]
\includegraphics[width=0.97\linewidth]{figures/benchgen_evaluations.png}
\caption{BenchGen evaluations. \fixme{POC: Sam -- this figure moves to Appendix.}}
\label{img:eval}
\label{evaluation_breakdown}
\end{figure}



Inspired by the recent paper (TBA: add references) on evaluation of image generation, we propose a graph representation to systematically represent prompts. Specifically, we use a tree to represent a major visual element appearing in a prompt. Within a tree, the root node represents the dominant object while child notes specify more details about the major element. Across trees, we introduce edges to represent the interaction between the trees, which can be a spatial relationship, action, or a dependency.

With this representation, we can systematically generate prompts while being able to analyze performances based on different criteria. For example, we can generate prompts with a certain number of trees or certain depths while also specifying how to expand the branches or interactions. For analysis, we can define the complexity based on the graph complexity and bucket prompts together by the number of trees or by the maximum number of dependencies.

Visual-Question-Answering-based Evaluations

For evaluation, this representation allows us to decompose a prompt into atomic visual elements. Measuring the model performance can then be just answering a series of simple questions. For example, to measure the text faithfulness of a generated image for the example prompt above, we can generate the following questions:
Is there a dog in the image?
Is the dog black?
Is the dog wearing glasses?
Are the glasses metal?
Is there an elephant in the picture?
Is the elephant happy?
Is there a great wall in the picture?
Is the dog taking a selfie with the elephant?
Are the dog and the elephant on the great wall?

By aggregating the answer to each of the questions, we can then compute the accuracy of the image with respect to this prompt. We can also further breakdown performance based on the type of questions we ask, \ie, is the question asking about interactions, or object generation?









\subsubsection{Annotation Criteria}
We use two separate evaluation metrics: visual appeal and text faithfulness. Visual appeal refers to the overall aesthetic quality of a generated image. It combines various visual elements such as color, shape, texture, and composition that creates a pleasing and engaging look. The concept of visual appeal is subjective and varies from person to person, as what may be aesthetically pleasing to one person may not be to another. Therefore, we ask five annotators to rate each sample. Concretely, we show annotators two images A and B, side-by-side, each generated by a different model using the same caption. The text captions are not displayed. Annotators choose which image is more visually appealing by selecting “A”, “B” or “Tie”. Text faithfulness refers to the degree of similarity between a generated image and a text caption. In this task, we again display two generated images A and B, side-byside, but with the caption alongside the images. The annotators are asked to ignore the visual appeal of the images and choose which ones best describe the caption with choices “A ”, “B”, “Both”, and “Neither”, where “Both” and “Neither” are considered as “Tie”. In this task, we have three annotators to annotate each sample pair. Wedonot report “standard” metrics such as FID scores. As argued in many recent papers, FID scores do not correlate well with human assessment of the performance of generative models.

We further expand to flaws for : human generation, text generation [TBA].


\subsubsection{Elo Ranking}


https://overleaf.thefacebook.com/project/66c3e87ee2df92c1c9667070



\subsection{\TextToV foundation evaluation}
\label{t2v_eval}
We evaluate \textToV foundation model on text-faithfulness, motion quality and video realness \& aesthetics.
Text-faithfulness evaluates the ability of the model aligns with prompt as instruction. Motion quality evaluates the ability of the model to understand object interactions in motion and generate consistent, natural motion on reasonable spatial parts. Realness and aesthetics evaluates the model's ability to generate close-to-real videos with good lighting/color/style, \etc.



\textbf{Motion Completeness}
We refer ``Motion Completeness'' as having enough amount of motion in generation. A lack of motion completeness happens when subject is out-of-distribution/unusual, for example, monsters, ghosts, or normal object doing unusual activity(people flying, panda playing piano). Because of lacking real world data, video generation model could have a harder time to generalize and generation appears to be more static or only contains camera motion.


For realness and aesthetic evaluation, we refer to realness as ``how close the video is to real video''. Realness as for out-of-distribution objects(monster/ghost) is defined as potentially from a movie shot. We refer to good aesthetic video as interesting content, cinematic lighting, descent color/saturation/exposure.








\subsubsection{Benchmark Generation}
1. Generate T2V

1)Performance \vs Gen3 (5s, 10s)
2)Performance \vs Luma (5s, 10s)
3)Performance \vs Kling (5s, 10s)
4)Performance \vs Sora(10s, 16s)

2. Storytelling

3. Precise Video Editing

\subsubsection{Annotation Criteria}
1. Video Quality Comparison

2. Prompt alignment Comparison

3.




\subsubsection{Elo Ranking}














